\begin{lemma}[Stabilization under the shift]
For every $n,k$,
\[
\mathsf{PowerCopiedLeft}_r(h,\mathsf{ShiftMap}(x),n,k)
=\mathsf{PowerCopiedLeft}_r(h,x,n+1,k).
\]
Thus the copied triangle constructed from the shifted input is exactly the
next row of the triangle constructed from the original input.
\end{lemma}
\begin{proof}
We prove the stronger assertion, simultaneously for every row $i$, that at
each depth $m$
\[
 \mathsf{PowerCopiedLeft}_r(h,\mathsf{ShiftMap}(x),i,m)
 =\mathsf{PowerCopiedLeft}_r(h,x,i+1,m).
 \tag{1}
\]
We induct on $m$.

For the base case $m=0$, fix an arbitrary row $i$.  The depth-zero clause
of \noderef{PowerCopiedLeft} gives
\[
\begin{aligned}
\mathsf{PowerCopiedLeft}_r(h,\mathsf{ShiftMap}(x),i,0)
 &=\mathsf{PowerIResponse}_r\bigl(
      h(\mathsf{PowerShiftIter}(i,\mathsf{ShiftMap}(x))),
      h(\mathsf{PowerShiftIter}(i+1,\mathsf{ShiftMap}(x)))\bigr),\\
\mathsf{PowerCopiedLeft}_r(h,x,i+1,0)
 &=\mathsf{PowerIResponse}_r\bigl(
      h(\mathsf{PowerShiftIter}(i+1,x)),
      h(\mathsf{PowerShiftIter}((i+1)+1,x))\bigr).
\end{aligned}
\tag{2}
\]
Apply \noderef{PowerShiftIterShift} first with index $i$ and then with
index $i+1$.  It yields
\[
\mathsf{PowerShiftIter}(i,\mathsf{ShiftMap}(x))
 =\mathsf{PowerShiftIter}(i+1,x)
\]
and
\[
\mathsf{PowerShiftIter}(i+1,\mathsf{ShiftMap}(x))
 =\mathsf{PowerShiftIter}((i+1)+1,x).
\]
Thus the first response in (2) and the second response in (2) receive the
same first argument and the same second argument.  Since
\noderef{PowerIResponse} is a fixed deterministic function of that ordered
pair of arguments, the two responses are equal.  This proves (1) at depth
zero for the arbitrary row $i$, and hence for every row.

For the successor step, assume (1) at depth $m$ for every row and fix an
arbitrary row $i$.  The induction hypothesis at row $i$ is
\[
\mathsf{PowerCopiedLeft}_r(h,\mathsf{ShiftMap}(x),i,m)
 =\mathsf{PowerCopiedLeft}_r(h,x,i+1,m),
\tag{3}
\]
while its specialization at row $i+1$ is
\[
\mathsf{PowerCopiedLeft}_r(h,\mathsf{ShiftMap}(x),i+1,m)
 =\mathsf{PowerCopiedLeft}_r(h,x,(i+1)+1,m).
\tag{4}
\]
The row arithmetic is $(i+1)+1=i+2$.  Therefore (3) and (4) identify the
two recursive values at rows $(i,i+1)$ on the shifted-input side with the
two recursive values at rows $(i+1,i+2)$ on the original-input side.

Unfold the successor-depth clause of \noderef{PowerCopiedLeft} on both
sides.  The left side at row $i$ and depth $m+1$ is
\[
\mathsf{PowerIResponse}_r\bigl(
 \mathsf{PowerCopiedLeft}_r(h,\mathsf{ShiftMap}(x),i,m),
 \mathsf{PowerCopiedLeft}_r(h,\mathsf{ShiftMap}(x),i+1,m)\bigr),
\tag{5}
\]
whereas the right side at row $i+1$ and depth $m+1$ is
\[
\mathsf{PowerIResponse}_r\bigl(
 \mathsf{PowerCopiedLeft}_r(h,x,i+1,m),
 \mathsf{PowerCopiedLeft}_r(h,x,(i+1)+1,m)\bigr).
\tag{6}
\]
By (3), the first inputs to the responses in (5) and (6) are equal; by
(4), the second inputs are equal.  The ordered input pairs are consequently
identical, and the determinism of \noderef{PowerIResponse} makes (5) and
(6) equal.  This proves (1) at depth $m+1$ for the arbitrary row $i$, so
the induction establishes (1) for every row and every depth.

Finally specialize (1) to $i=n$ and $m=k$.  This is exactly the equality
claimed in the lemma.
\end{proof}
